\subsection{Transpose of a continuous linear mapping}

In this section, \(E_{1}\) and \(E_{2}\) denote two locally convex spaces, with respective duals \(E_{1}^{\prime}\) and \(E_{2}^{\prime}\) .

Let \(u\) be a linear mapping from \(\mathrm{E}_1\) into \(\mathrm{E}_2\) . For \(u\) to be continuous when \(\mathrm{E}_1\) and \(\mathrm{E}_2\) are assigned the weakened topologies, it is necessary and sufficient that \(f \circ u\) belongs to \(\mathrm{E}_1'\) for all \(f \in \mathrm{E}_2'\) ; this is the case if \(u\) is continuous. Then the linear mapping \(f \mapsto f \circ u\) from \(\mathrm{E}_2'\) into \(\mathrm{E}_1'\) is called the transpose of \(u\) and is denoted by \(^t u\) .

PROPOSITION 5. --- Let u be a continuous linear mapping from \(E_{1}\) into \(E_{2}\) .

\begin{enumerate}
\def\labelenumi{(\roman{enumi})}
\item
  If \(E_{1}\) and \(E_{2}\) are Hausdorff then u is injective if and only if the image of \(^{t}u\) is dense in \(E_{1}^{\prime}\) for the weak topology \(\sigma(\mathrm{E}_{1}^{\prime},\mathrm{E}_{1})\) .
\item
  For \({}^{t}u\) to be injective it is necessary and sufficient that the image of \(u\) is dense in \(\mathbf{E}_2\) .
\end{enumerate}

A vector subspace of \(E_{2}\) is dense for the initial topology if and only if it is dense for the weakened topology (IV, p.~4, prop. 2). Prop. 5 then follows from II, p.~47, cor. 2.

PROPOSITION 6. --- Let u be a linear mapping from \(E_{1}\) into \(E_{2}\) which is continuous for the weakened topologies. For i = 1, 2, let \(S_{i}\) be a family of bounded subsets of \(E_{i}\) . In order that \(^{t}u\) is a continuous mapping from \((\mathrm{E}_{2}^{\prime})_{\mathfrak{S}_{2}}\) into \((\mathrm{E}_{1}^{\prime})_{\mathfrak{S}_{1}}\) , it is necessary and sufficient that, for every set \(A \in S_{1}\) , there exist sets \(A_{1}, \ldots, A_{n}\) in \(S_{2}\) and a real number \(\lambda > 0\) such that \(\lambda.u(\mathbf{A})\) is contained in the closed convex balanced envelope of \(A_{1} \cup \ldots \cup A_{n}\) .

This is an immediate consequence of prop. 2 of III, p.~15.

COROLLARY. --- Let u be a continuous linear mapping from \(E_{1}\) into \(E_{2}\) . Then \({}^t u\) is continuous when the duals \(E_{i}'\) are assigned the following topologies :

\begin{enumerate}
\def\labelenumi{\alph{enumi})}
\item
  the weak topologies \(\sigma (\mathrm{E}_i^{\prime},\mathrm{E}_i)\)
\item
  the strong topologies \(\beta (\mathrm{E}_i^{\prime},\mathrm{E}_i)\)
\item
  the Mackey topologies \(\tau (\mathrm{E}_i^{\prime},\mathrm{E}_i)\)
\item
  the topologies of precompact convergence.
\end{enumerate}

Moreover, if \(E_{2}\) is Hausdorff, \(^{t}u\) is continuous when the duals \(E_{i}^{\prime}\) are assigned:

\begin{enumerate}
\def\labelenumi{\alph{enumi})}
\setcounter{enumi}{4}
\tightlist
\item
  the topologies of compact convergence (resp. compact convex).
\end{enumerate}

The only point which requires a proof is the case \(c\) ), when the topologies of \(\mathrm{E}_1\) and \(\mathrm{E}_2\) are not necessarily Hausdorff. Then for every linear form \(f \in \mathrm{E}'_1{}^*\) , \(f \circ {}^t u\) is a linear form on \(\mathrm{E}_2'\) ; hence there is a linear mapping \(v: \mathrm{E}'_1{}^* \to \mathrm{E}'_2{}^*\) which is continuous for the topologies \(\sigma(\mathrm{E}'_1{}^*, \mathrm{E}'_1)\) and \(\sigma(\mathrm{E}'_2{}^*, \mathrm{E}'_2)\) and is such that \(d_{\mathrm{B}_2} \circ u = v \circ d_{\mathrm{B}_1}\) , where \(d_{\mathrm{B}_i}\) is the canonical mapping from \(\mathrm{E}_i\) into \(\mathrm{E}'_i{}^* (i = 1, 2)\) . Consequently, if A is a subset of \(\mathrm{E}_1\) such that \(d_{\mathrm{B}_1}(\mathrm{A})\) is convex, balanced and compact for \(\sigma(\mathrm{E}'_1{}^*, \mathrm{E}'_1)\) then \(d_{\mathrm{B}_2}(u(\mathrm{A})) = v(d_{\mathrm{B}_1}(\mathrm{A}))\) is convex, balanced and compact for \(\sigma(\mathrm{E}'_2{}^*, \mathrm{E}'_2)\) since the topologies \(\sigma(\mathrm{E}'_1{}^*, \mathrm{E}'_1)\) and \(\sigma(\mathrm{E}'_2{}^*, \mathrm{E}'_2)\) are Hausdorff.

PROPOSITION 7. --- Let \(u: \mathrm{E}_1 \to \mathrm{E}_2\) be a linear mapping. We assume that \(u\) is continuous for the weakened topologies of \(\mathrm{E}_1\) and \(\mathrm{E}_2\) .

\begin{enumerate}
\def\labelenumi{(\roman{enumi})}
\item
  The mapping \(u\) is continuous if \(\mathrm{E}_1\) and \(\mathrm{E}_2\) are assigned their Mackey topologies.
\item
  If \(\mathbf{E}_1\) is bornological or barrelled, then \(u\) is continuous for the initial topologies of \(\mathbf{E}_1\) and \(\mathbf{E}_2\) .
\item
  In order that \(u\) be continuous for the initial topologies of \(\mathbf{E}_1\) and \(\mathbf{E}_2\) , it is necessary and sufficient that the image under \(^t u\) of every equicontinuous subset of \(\mathbf{E}_2'\) be equicontinuous in \(\mathbf{E}_1'\) .
\end{enumerate}

The hypothesis implies that \({}^t u\) is continuous for the weak topologies \(\sigma(\mathrm{E}_2', \mathrm{E}_2)\) and \(\sigma(\mathrm{E}_1', \mathrm{E}_1)\) (II, p.~46, corollary) hence the image under \({}^t u\) of a convex, balanced and compact subset for \(\sigma(\mathrm{E}_2', \mathrm{E}_2)\) is convex, balanced and compact for \(\sigma(\mathrm{E}_1', \mathrm{E}_1)\) , the topologies \(\sigma(\mathrm{E}_2', \mathrm{E}_2)\) and \(\sigma(\mathrm{E}_1', \mathrm{E}_1)\) being Hausdorff. Therefore, assertion (i) follows from GT, X, § 1, No.~4, prop. 3, b). Assertion (ii) is a consequence of (i): for, if \(\mathrm{E}_1\) is bornological or barrelled, its initial topology is the Mackey topology, and the Mackey topology of \(\mathrm{E}_2\) is finer than the initial topology of \(\mathrm{E}_2\) . Finally, the initial topology of \(\mathrm{E}_i\) is that of uniform convergence on equicontinuous subsets of \(\mathrm{E}_i'\) (III, p.~19, cor. 1 of prop. 7). This proves (iii).

COROLLARY. --- Suppose \(E_{1}\) is a normed space. Let u be a linear mapping from \(E_{1}\) into \(E_{2}\) . The following properties are equivalent :

\begin{enumerate}
\def\labelenumi{\alph{enumi})}
\item
  \(u\) is continuous;
\item
  \(u\) is continuous for the weakened topologies;
\item
  the image of the unit ball in \(\mathrm{E}_1\) under \(u\) is bounded in \(\mathrm{E}_2\) ;
\item
  for every sequence \((x_{n})\) of points of \(\mathbf{E}_1\) tending to 0 for the initial topology, the sequence \((u(x_{n}))\) is bounded for the weakened topology of \(\mathbf{E}_2\) .
\end{enumerate}

Since \(E_{1}\) is bornological the equivalence of a) and b) follows from prop. 7; that of a) and c) is immediate. The equivalence of a) and d) follows from prop. 1 of IV, p.~1, and from prop. 1 of III, p.~11.

PROPOSITION 8. --- (i) Let E be a normed space, with dual E'. For every \(x \in E\) , we have

\[
\| x \| = \sup _ {x ^ {\prime} \in \mathrm{E} ^ {\prime}, \| x ^ {\prime} \| \leqslant 1} | \langle x, x ^ {\prime} \rangle |.\tag{3}
\]

\begin{enumerate}
\def\labelenumi{(\roman{enumi})}
\setcounter{enumi}{1}
\tightlist
\item
  Let \(\mathbf{E}_1\) and \(\mathbf{E}_2\) be two normed spaces and \(u\) a continuous linear mapping from \(\mathbf{E}_1\) into \(\mathbf{E}_2\) . We have
\end{enumerate}

\[
\left\| ^ {t} u \right\| = \left\| u \right\|.\tag{4}
\]

Let \(x \in E\) . For every \(x' \in E'\) such that \(\| x' \| \leqslant 1\) , we have

\[
\left| \langle x, x ^ {\prime} \rangle \right| \leqslant \| x \| \| x ^ {\prime} \| \leqslant \| x \|.
\]

By Hahn-Banach theorem (II, p.~23, cor. 2), there exists an element \(x'\) in \(E'\) such that \(\| x' \| \leqslant 1\) and \(\langle x, x' \rangle = \| x \|\) . This proves (i).

Let us now prove (ii). By formula (3) and the definition of the transpose, we have

\[
\begin{array}{l} \| ^ {t} u \| = \sup _ {\| y ^ {\prime} \| \leqslant 1} \left\| ^ {t} u (y ^ {\prime}) \right\| = \sup _ {\| y ^ {\prime} \| \leqslant 1, \| x \| \leqslant 1} \left| \langle x, ^ {t} u (y ^ {\prime}) \rangle \right| \\ = \sup _ {\| x \| \leqslant 1, \| y ^ {\prime} \| \leqslant 1} \left| \langle u (x), y ^ {\prime} \rangle \right| = \sup _ {\| x \| \leqslant 1} \left\| u (x) \right\| = \| u \|. \end{array}
\]

Remarks. --- 1) Formula (3) is a particular case of (4), corresponding to the linear mapping \(\lambda \mapsto \lambda x\) from K into E.

\begin{enumerate}
\def\labelenumi{\arabic{enumi})}
\setcounter{enumi}{1}
\tightlist
\item
  Put \(\mathrm{B}(x, y') = \langle u(x), y' \rangle = \langle x, {}^t u(y') \rangle\) for \(x \in \mathrm{E}_1\) , \(y' \in \mathrm{E}_2'\) . The above proof shows that \(\mathbf{B}\) is a continuous bilinear form on \(\mathrm{E}_1 \times \mathrm{E}_2'\) , with norm (GT, X, § 3, No.~2) equal to \(\| u \|\) .
\end{enumerate}

COROLLARY. --- Let E be a normed space satisfying the first axiom of countability. There exists a countable subset D of \(E' - \{0\}\) such that we have

\[
\| x \| = \sup _ {\xi \in \mathbf {D}} | \langle x, \xi \rangle | / \| \xi \|\tag{5}
\]

for all \(x \in E\) .

Let \(B'\) be the unit ball of the dual \(E'\) of \(E\) with the weak topology \(\sigma(E', E)\) assigned to it. Then \(B'\) is a compact metrizable space (III, p.~19, cor. 2); hence there exists a countable dense subset \(D'\) in \(B'\) . Put \(D = D' \cap (E' - \{0\})\) . Let \(x \in E\) ; the mapping \(x' \mapsto \langle x, x' \rangle\) from \(B'\) into \(K\) is continuous, therefore

\[
\sup _ {x ^ {\prime} \in \mathbf {B} ^ {\prime}} | \langle x, x ^ {\prime} \rangle | = \sup _ {\xi \in \mathbf {D} ^ {\prime}} | \langle x, \xi \rangle | \leqslant \sup _ {\xi \in \mathbf {D}} | \langle x, \xi \rangle | / \| \xi \| \leqslant \| x \|.
\]

Formula (5) now follows from (3).

